
\medskip
\noindent \SGARefTarget{sga7.sga7i.eviii.n2-4}{2.4}.  Let
\[
 f:\underline T'\longrightarrow\underline T
\]
be a morphism of topoi.  There is then a canonical homomorphism
(VII \SGARefLink{sga7.sga7i.evii.n2-8}{2.8})
\[
 \Biext^1(P,Q;G)\longrightarrow\Biext^1(P',Q';G'),
\]
where primes denote inverse-image groups.  Since, for every
commutative group $E$ of $\underline T$, there is plainly a canonical
isomorphism
\[
 ({}_nE)'\simeq{}_n(E'),
\]
the homomorphisms (\SGARefLink{sga7.sga7i.eviii.n2-2-2}{2.2.2}) over $\underline T$ and $\underline T'$ give
a canonical diagram
\begin{equation}\tag{\SGARefTarget{sga7.sga7i.eviii.n2-4-1}{2.4.1}}
\begin{tikzcd}[column sep=huge,row sep=huge]
 \Biext^1(P,Q;G)\arrow[r]\arrow[d,"{\varphi^n_{P,Q}}"']
   &\Biext^1(P',Q';G')\arrow[d,"{\varphi^n_{P',Q'}}"]\\
 \Hom({}_nP\otimes{}_nQ,G)\arrow[r]
   &\Hom({}_nP'\otimes{}_nQ',G').
\end{tikzcd}
\end{equation}
Similarly, the homomorphisms (\SGARefLink{sga7.sga7i.eviii.n2-2-3}{2.2.3}) give an analogous diagram in
terms of the $T_\ell$.  I claim that these diagrams are
\underline{commutative}.  The verification is essentially immediate
from the definitions and is left to the reader.  More generally,
notice that \underline{all} the canonical homomorphisms considered in
this section and the preceding one are \underline{compatible with
change of topos}, in the evident sense.  We shall henceforth use
compatibilities of this kind without further mention.

